\label{ch:modular-pva-quantization}%
\label{def:complete-modular-bar}%

\begin{abstract}
This note reorganizes the entire discussion of the present thread into a single fortified
package, written from first principles and separated carefully into three layers:

\begin{enumerate}[label=\textnormal{(\roman*)},leftmargin=2.6em]
\item \emph{external input} from the literature;
\item \emph{formal algebraic results} proved internally in this note once a modular bar datum is given;
\item \emph{conditional comparison statements} that connect the algebraic package to the
three-dimensional holomorphic-topological Poisson sigma model of Khan--Zeng.
\end{enumerate}

The central principle is that \emph{trees govern chirality and loops govern modularity}.
Disk-level quadratic duality controls only genus zero. A genuinely modular deformation
theory must already contain a primitive genus-raising one-edge operator. We make this
precise by constructing a stable-graph modular bar coalgebra, defining the complete
filtered strict modular deformation dg Lie algebra
\[
\Defmod(A_{\cl}(V)) := \Coder(\Bmod(\A))[-1],
\]
and proving the universal obstruction recursion for lifting a genus-zero Maurer--Cartan
class to all genera. We then pass to a local variational coordinate model for filtered
quadratic Poisson vertex algebras, identify the first genus-raising operator with a reduced
odd Laplacian
\[
D_1=\Delta_{\mathrm{cyc}},
\]
construct the associated variational modular class, prove a vanishing theorem for a broad
triangular $W$-normal form, and compute explicitly that the first modular obstruction
vanishes for classical $W_3$.

Finally, in the local, translation-invariant, filtration-preserving, reduced conformal-weight
zero sector, we compute the relevant shifted $H^1$ for $W_3$ and prove that the only
nontrivial genus-one lift direction is the central-parameter direction $\partial_c P_c$.
Thus, in the relevant one-loop sector, every genus-one lift is gauge equivalent to a
renormalization of the central parameter.

The resulting picture gives a precise candidate for the missing modular homotopy theory
behind deformation quantization of Poisson vertex algebras to vertex algebras, while
remaining honest about which inputs are presently formal, conditional, or conjectural.
\end{abstract}

\section{Status conventions and scope}

The note is intentionally stratified.

\begin{convention}[Three levels of assertion]
Throughout the note we distinguish three kinds of statements.
\begin{enumerate}[leftmargin=2.4em]
\item \textbf{External input.} These are statements imported from the literature:
Khan--Zeng on the three-dimensional holomorphic-topological Poisson sigma model,
Gui--Li--Zeng on quadratic duality for chiral algebras, Francis--Gaitsgory on chiral
Koszul duality, Costello--Dimofte--Gaiotto on boundary chiral algebras in 3d
holomorphic twists, Gaiotto--Kulp--Wu on higher operations, and related works
\cite{KZ25,GLZ22,FG12,CDG20,GKW25}.
\item \textbf{Formal algebraic results.} These are theorems proved in this note once one
has fixed a modular bar datum. They include the stable-graph bar construction, the
filtered strict modular deformation dg Lie algebra, the universal obstruction recursion, the coordinate
formula for the genus-raising operator, the variational modular class, the triangular
$W$-normal-form vanishing theorem, and the $W_3$ relevant cohomology computation.
\item \textbf{Conditional comparison results.} These are statements linking the formal
algebra to the Poisson sigma model. They depend on analytic and geometric inputs
not proved here, such as the existence of the appropriate resolved chiral/factorization
datum attached to a given PVA and the existence of half-space perturbative quantization.
\end{enumerate}
\end{convention}

\begin{remark}[Why this stratification matters]
The algebraic stable-graph
package below is self-contained once a modular bar datum is fixed, but the extraction of
that datum from a general PVA via compactified configuration spaces and the identification
with the three-dimensional Poisson sigma model remain conditional at present.
\end{remark}

\begin{convention}[Strict models versus invariant modular deformation objects]
\label{conv:modular-pva-strict-versus-homotopy}
\index{modular deformation!strict versus homotopy}
Throughout this note, concrete coderivation dg~Lie algebras such as
$\Defmod(A_{\cl}(V))$ are chosen strict presentations.  The invariant
modular deformation theory is defined only up to filtered
$L_\infty$ quasi-isomorphism.  Accordingly, every Maurer--Cartan equation
written in the strict modular complex should be read as a computation in
one preferred chart on an underlying homotopy-coherent modular
deformation object $\Definfmod(\cA)$ constructed in Volume~I.
\end{convention}

\section{External input from the literature}

We record the precise ingredients from the literature that motivate and support the
new constructions.

\subsection{Poisson vertex algebras and the 3d holomorphic-topological Poisson sigma model}

Khan and Zeng introduce a mixed holomorphic-topological gauge theory in three dimensions
attached to a freely generated Poisson vertex algebra (PVA), show that gauge invariance
holds if and only if the $\lambda$-bracket Jacobi identity holds, prove that the presence
of a Virasoro element upgrades the symmetry from holomorphic-topological to fully
topological, and propose the resulting three-dimensional Poisson sigma model as a route
toward deformation quantization of PVAs to vertex algebras \cite{KZ25}. Their
boundary picture on the half-space $\mathbb R_+\times \mathbb C$ is the physical source of
the genus-zero quantization problem considered here.

\subsection{Quadratic duality for chiral algebras}

Gui, Li, and Zeng define quadratic duality for chiral algebras, prove an injective map
\[
\Hom(A,B)\hookrightarrow \MC(A^!\otimes B),
\]
show bijectivity under effectiveness hypotheses in degree zero, and in the inhomogeneous
quadratic setting introduce twisted pairs $(B,B^\circ,S)$ together with the curved
Maurer--Cartan equation
\[
\mu\big((S+\alpha)\boxtimes (S+\alpha)\big)=0.
\]
They also state explicitly that a full homotopy/Koszul treatment, parallel to the
associative story, is still missing \cite{GLZ22}. This gap is one of the main
algebraic motivations for the present note.

\subsection{Chiral/factorization homotopy theory}

Francis and Gaitsgory place chiral and factorization structures in a common homotopy-theoretic
framework and recast their relation as a form of chiral Koszul duality
\cite{FG12}. This supplies the correct derived ambient category in which
one should expect a modular refinement of disk-level chiral duality to live.

\subsection{Boundary algebras, higher operations, and line operators}

Costello, Dimofte, and Gaiotto show that in 3d holomorphic twists bulk local operators form
a commutative chiral algebra with shifted Poisson bracket, while boundary operator algebras
naturally carry higher $A_\infty$ structure \cite{CDG20}. Gaiotto,
Kulp, and Wu construct higher operations in holomorphic-topological perturbation theory,
prove the relevant quadratic consistency identities, and interpret them as a Wess--Zumino
type consistency for BRST symmetry \cite{GKW25}. Dimofte, Niu, and Py show
that line operators in perturbative 3d holomorphic QFT are represented by modules for a
Koszul-dual $A_\infty$-algebra $A^!$ and are organized by Maurer--Cartan data with an
$A_\infty$ Yang--Baxter flavor \cite{DNP25}. All of this reinforces the thesis
that \emph{a single Maurer--Cartan object should package higher operations, anomalies,
and OPE compatibilities}. 

\subsection{Analogies with deformation quantization}

Cattaneo and Felder interpret Kontsevich's deformation quantization via a topological
sigma model on the disk \cite{CF00}. Tamarkin proves that the deformation
complex of an $E_d$-algebra carries an $E_{d+1}$-structure \cite{Tam02}. These
results suggest that the correct receptacle for quantization of PVAs should not be merely
a disk-level dg Lie algebra, but a higher and genuinely modular deformation object.

\section{From first principles: genus zero, loops, and the one-edge principle}

The main structural observation can be stated in one sentence:

\begin{center}
\emph{A deformation theory that only remembers trees cannot see modular obstructions.}
\end{center}

We make that precise.

\subsection{Filtered quadratic PVAs}

\begin{definition}[Filtered quadratic PVA]
A \emph{filtered quadratic Poisson vertex algebra} is a differential commutative algebra
freely generated by fields $u^1,\dots,u^N$ and their derivatives, together with a
$\lambda$-bracket encoded by a skew-adjoint matrix differential operator
\[
\Pi^{ab}(\partial)=\sum_{m\ge 0}\Pi^{ab}_m\,\partial^m,
\]
whose coefficients are filtered polynomials of total generator degree at most $2$.
\end{definition}

The classical $W$-algebras discussed by Khan--Zeng fit this framework, including the
classical $W_3$ example \cite{KZ25}.

\subsection{Resolved classical datum}

\begin{assumption}[Resolved classical chiral/factorization datum]
Let $X$ be a smooth complex curve and let $V$ be a filtered quadratic PVA on $X$.
Assume there exists a cofibrant resolved classical chiral/factorization object
\[
\A=\widetilde A_{\cl}(V)
\]
whose genus-zero shadow is $V$, and which is dualizable, finite type, conilpotent,
and complete in the sense made precise below.
\end{assumption}

\begin{remark}
This assumption isolates the step that is currently external to the purely algebraic theory.
The remainder of the formal constructions depend only on the existence of an input
modular bar datum built from $\A$.
\end{remark}

\subsection{The one-edge principle}

Taking the disk-level deformation complex and adding a genus grading does not suffice.

\begin{proposition}[Necessity of a genus-raising primitive; \ClaimStatusProvedHere]
Suppose a filtered deformation dg Lie algebra $L$ is equipped with a complete genus
filtration $F^\bullet$ and its differential preserves genus exactly. Then any genus-zero
Maurer--Cartan element $\Theta_0\in \MC(\gr_F^0 L)$ extends tautologically to a full
Maurer--Cartan element of $L$ without further obstruction.
\end{proposition}

\begin{proof}
If the differential preserves genus and the Lie bracket adds genera, then the Maurer--Cartan
equation decomposes strictly by genus. A solution in genus zero has no forcing term in
positive genus, so the same element, viewed inside the completed filtered algebra, already
solves the full equation.
\end{proof}

\begin{remark}[Trees versus loops]
Therefore any nontrivial modular obstruction theory must contain, at the chain level,
a primitive operator that \emph{raises genus}. In the stable-graph bar construction below,
that primitive is the nonseparating one-edge expansion operator. This is the algebraic
counterpart of the first loop correction in perturbation theory.
\end{remark}

\section{Stable-graph modular bar data and the modular bar coalgebra}

This section is the formal heart of the new algebraic package.

\subsection{Stable graphs}

\begin{definition}[Connected stable graph]
A \emph{connected stable graph} $\Gamma$ consists of the following data:
\begin{enumerate}[leftmargin=2.4em]
\item a finite set $V(\Gamma)$ of vertices;
\item a finite set $\mathrm{Fl}(\Gamma)$ of flags;
\item a partition $\mathrm{Fl}(\Gamma)=\bigsqcup_{v\in V(\Gamma)}\mathrm{Fl}(v)$;
\item an involution on $\mathrm{Fl}(\Gamma)$, whose $2$-cycles are internal edges and
whose fixed points are external legs;
\item a genus label $g(v)\in \mathbb Z_{\ge 0}$ for each vertex,
\end{enumerate}
subject to the stability condition
\[
2g(v)-2+|\mathrm{Fl}(v)|>0
\qquad\text{for all }v\in V(\Gamma).
\]
Its total genus is
\[
g(\Gamma)=b_1(\Gamma)+\sum_{v\in V(\Gamma)}g(v).
\]
\end{definition}

Write $E(\Gamma)$ for the set of internal edges and
\[
\orline(\Gamma):=\det(kE(\Gamma))\otimes \det(H_1(\Gamma;k))^{-1}
\]
for the orientation line.

\subsection{One-edge expansions}

Let $C$ be a reduced complete finite-type graded collection on finite sets:
\[
F\longmapsto C(F),
\]
functorial for bijections of finite sets.

\begin{definition}[One-edge expansion set]
For a finite set $F$, let $\OE(F)$ denote the set of connected stable graphs whose
external leg set is $F$ and which have exactly one internal edge. Such a graph is either:
\begin{enumerate}[leftmargin=2.4em]
\item \emph{separating}: two vertices joined by one edge;
\item \emph{nonseparating}: one vertex carrying one loop.
\end{enumerate}
\end{definition}

\begin{definition}[Modular bar datum]
\label{def:modular-bar-datum}
A \emph{modular bar datum} on $C$ consists of:
\begin{enumerate}[leftmargin=2.4em]
\item an internal differential
\[
d:C(F)\to C(F)
\]
of degree $+1$ for each finite set $F$;
\item for each $\epsilon\in \OE(F)$, a degree $+1$ map
\[
\Delta_\epsilon:
s\,C(F)\longrightarrow
\orline(\epsilon)\otimes \bigotimes_{u\in V(\epsilon)} s\,C(\mathrm{Fl}_\epsilon(u));
\]
\item the \emph{codimension-two cancellation axiom}: for every connected stable graph
$T$ with exactly two internal edges, the signed sum of the two ways of producing $T$
by successive one-edge expansions vanishes;
\item the anticommutation relation that $d$ anticommutes with every $\Delta_\epsilon$.
\end{enumerate}
\end{definition}

\begin{remark}[Why this is the right primitive datum]
The usual bar construction is governed by codimension-one collisions on configuration
spaces. In the modular setting, codimension-one boundary strata correspond to one-edge
degenerations of stable graphs. Thus the primitive input is not a generic higher operator
but precisely the family of one-edge expansion maps.
\end{remark}

\subsection{The complete modular bar coalgebra}
\label{sec:modular-bar-coalgebra}

For a connected stable graph $\Gamma$, define
\[
C[\Gamma]:=
\orline(\Gamma)\otimes\bigotimes_{v\in V(\Gamma)} s\,C(\mathrm{Fl}(v)).
\]

\begin{definition}[Complete modular bar object]
The \emph{complete modular bar object} of a modular bar datum $C$ is
\[
\Bmod(C):=
\prod_{[\Gamma]\in \StGraph^{\mathrm{conn}}}
\big(C[\Gamma]\big)_{\mathrm{Aut}(\Gamma)},
\]
completed with respect to total genus.
\end{definition}

\begin{definition}[Internal and expansion differentials]
Define degree-$+1$ endomorphisms
\[
D_{\mathrm{int}},\ D_{\mathrm{exp}}:\Bmod(C)\to \Bmod(C)
\]
on graph generators as follows.
\begin{enumerate}[leftmargin=2.4em]
\item $D_{\mathrm{int}}$ applies the internal differential $d$ to one vertex label and
sums over vertices.
\item $D_{\mathrm{exp}}$ replaces a chosen vertex $v$ by a one-edge expansion
$\epsilon\in \OE(\mathrm{Fl}(v))$ and applies $\Delta_\epsilon$ to the label at $v$,
with the sign determined by the orientation line and Koszul rule.
\end{enumerate}
Set
\[
D:=D_{\mathrm{int}}+D_{\mathrm{exp}}.
\]
\end{definition}

\begin{theorem}[Stable-graph modular bar construction; \ClaimStatusProvedHere]\label{thm:modular-bar}
Let $C$ be a modular bar datum. Then:
\begin{enumerate}[leftmargin=2.4em]
\item $D^2=0$;
\item $\Bmod(C)$ is a complete conilpotent modular dg coalgebra;
\item if $F^g\Bmod(C)$ denotes the subspace spanned by graphs of total genus $\ge g$, then
\[
\Bmod(C)\cong \varprojlim_g \Bmod(C)/F^g\Bmod(C),
\]
and this filtration is exhaustive and complete.
\end{enumerate}
\end{theorem}

\begin{proof}
The square of $D$ has three types of terms.

\smallskip
\noindent
\emph{(i) Two internal differentials.}
These vanish because $d^2=0$ on every vertex label.

\smallskip
\noindent
\emph{(ii) One internal differential and one expansion.}
These cancel because $d$ anticommutes with every one-edge expansion map $\Delta_\epsilon$.

\smallskip
\noindent
\emph{(iii) Two one-edge expansions.}
These are exactly the codimension-two stable-graph boundary terms. Every graph with two
internal edges can be obtained by two successive one-edge expansions, and the codimension-two
cancellation axiom says the signed sum of the two resulting paths is zero.

Therefore $D^2=0$.

For conilpotence, use the reduced coproduct defined by extracting a nontrivial connected
stable subgraph. Repeated reduced coproduct strictly decreases the number of internal
edges, hence eventually vanishes on every graph class. Completeness is immediate from
the completed product over total genus.
\end{proof}

\begin{remark}[Tree-level truncation]
The genus-zero subquotient of $\Bmod(C)$ is obtained by restricting to connected stable
graphs of total genus zero, i.e.\ trees all of whose vertices have genus zero. We denote
that tree-level coalgebra by $\Bch(C)$.
\end{remark}

\subsection{Separating and nonseparating pieces}

Split the expansion operator into
\[
D_{\mathrm{exp}}=D_{\mathrm{sep}}+D_{\mathrm{nsep}},
\]
where $D_{\mathrm{sep}}$ sums over separating one-edge expansions and $D_{\mathrm{nsep}}$
over nonseparating one-edge expansions. Set
\[
D_0:=D_{\mathrm{int}}+D_{\mathrm{sep}},
\qquad
D_1:=D_{\mathrm{nsep}}.
\]

\begin{proposition}[The one-edge decomposition; \ClaimStatusProvedHere]\label{prop:D0D1}
With the notation above:
\begin{enumerate}[leftmargin=2.4em]
\item $D_0$ preserves total genus;
\item $D_1$ raises total genus by one;
\item $D_0^2=0$, $D_0D_1+D_1D_0=0$, and $D_1^2=0$.
\end{enumerate}
\end{proposition}

\begin{proof}
Separating one-edge expansions preserve the first Betti number and total genus, while
nonseparating one-edge expansions increase the first Betti number by one. The identities
follow by taking the genus-homogeneous pieces of the equality $(D_0+D_1)^2=0$.
\end{proof}

\subsection{Curved extension via $S$-tails}

Gui--Li--Zeng show that in the inhomogeneous quadratic setting the correct input is not
merely a CDG algebra, but a twisted pair $(B,B^\circ,S)$, and that the passage to the
associated CDG algebra forgets information \cite{GLZ22}. The modular extension
should therefore remember the element $S$ geometrically.

\begin{definition}[Curved one-tail expansion]
Given a twisted pair $(B,B^\circ,S)$, an \emph{$S$-tail expansion} is the operation that
attaches a new univalent $S$-decorated vertex to a pre-existing corolla. Formally, it is
a unary expansion map
\[
\Delta_S:s\,C(F)\to s\,C(F\sqcup \{\bullet,\star\})
\]
followed by contraction of the new internal edge, with the new univalent vertex labeled by $S$.
\end{definition}

\begin{remark}
The correct curved modular bar object is obtained by adjoining the $S$-tail expansions
to the set of primitive degenerations. Provided they satisfy the same codimension-two
compatibilities with the one-edge expansions, the proof of \cref{thm:modular-bar} goes
through verbatim. This is the modular analogue of the twisted Maurer--Cartan equation
\[
\mu\big((S+\alpha)\boxtimes (S+\alpha)\big)=0
\]
of Gui--Li--Zeng.
\end{remark}

\section{The strict modular deformation dg Lie algebra}

We now pass from the coalgebra to the deformation complex.

\begin{definition}[Modular deformation algebra]
Let $C$ be a modular bar datum. Define
\[
L_{\mathrm{mod}}(C):=\Coder(\Bmod(C))[-1].
\]
If $\A$ is a resolved classical chiral/factorization datum with underlying modular bar datum
$C(\A)$, write
\[
\Defmod(A_{\cl}(V)):=L_{\mathrm{mod}}(C(\A)).
\]
\end{definition}

\begin{remark}[Strict model of the modular deformation object]
\label{rem:lmod-strict-model}
The dg~Lie algebra $L_{\mathrm{mod}}(C)$ is the strict model of
the homotopy-invariant modular deformation object
$\Definfmod(\cA)$, by the modular homotopy-convolution theorem of
Volume~I.
This chapter's filtered obstruction theory
(Theorems~\ref{thm:Ob1},~\ref{thm:Obg}) operates at the strict
level because the complete filtration ensures convergence of the
MC tower.
\end{remark}

\begin{proposition}[Convolution description; \ClaimStatusProvedHere]
A coderivation of $\Bmod(C)$ is determined by its projection to corollas. Hence there is
a canonical identification
\[
L_{\mathrm{mod}}(C)\cong \Hom(\Bmod(C),sC)[-1]
\]
with dg Lie bracket given by commutator.
\end{proposition}

\begin{proof}
This is the standard universal property of cofree conilpotent coalgebras, applied in the
completed stable-graph setting.
\end{proof}

\begin{construction}[The homotopy modular deformation object]
\label{const:vol2-homotopy-modular-object}
\index{modular deformation!homotopy object}
Let $C$ be a modular bar datum.  Denote by
\[
\Definfmod(C)
\]
the complete filtered $L_\infty$-algebra obtained from any cofibrant or
quasi-free resolution of the stable-graph bar coalgebra, or equivalently
from the coderivation complex of a bar-cofree strictification.  Its
Maurer--Cartan elements are the homotopy-coherent modular twisting data
attached to~$C$, and any two such constructions are related by filtered
$L_\infty$ quasi-isomorphisms, exactly as in the modular homotopy
deformation-object construction of Volume~I.
\end{construction}


\subsection{From trees to stable graphs: the genus completion}
\label{subsec:genus-completion}
\index{genus completion|textbf}

The modular bar coalgebra $\Bmod(C)$ carries a complete
genus filtration.  Its genus-$0$ piece is the ordinary
tree-level bar coalgebra; the modular deformation algebra
$L_{\mathrm{mod}}(C)$ is the genus completion of the
tree-level deformation complex.  The non-separating one-edge
expansion is the sole primitive that connects them.

\begin{definition}[Chiral convolution dg~Lie algebra]
\label{def:chiral-convolution}
\index{chiral convolution dg Lie algebra|textbf}
The \emph{tree-level} or \emph{chiral convolution dg~Lie algebra}
of a modular bar datum $C$ is
\begin{equation}\label{eq:chiral-convolution-dg-lie}
\mathfrak{g}^{\mathrm{ch}}_C
\;:=\;
\Coder(\Bch(C))[-1]
\;\cong\;
\Hom(\Bch(C),\, sC)[-1],
\end{equation}
where $\Bch(C) := \bar{B}^{(0)}(C)$ is the tree-level
sub-coalgebra of~$\Bmod(C)$---stable graphs with all vertex
genera zero and $b_1 = 0$, i.e.\ trees.  Its MC elements are
tree-level twisting cochains.
\end{definition}

\begin{theorem}[Genus completion; \ClaimStatusProvedHere]
\label{thm:genus-completion}
\index{genus filtration!on modular deformation algebra}
\index{modular deformation algebra!genus completion}
The genus filtration
$F^g\Bmod(C) := \prod_{g(\Gamma) \ge g} C[\Gamma]_{\Aut(\Gamma)}$
induces a complete descending filtration
$F^\bullet L_{\mathrm{mod}}(C)$ with the following properties:
\begin{enumerate}[label=\textup{(\roman*)}]
\item The genus-$g$ associated graded of the coalgebra is the
  genus-$g$ bar complex:
  \[
  \gr_F^g \Bmod(C)
  \;\cong\;
  \bar{B}^{(g)}(C)
  \;=\;
  \bigoplus_{\substack{\Gamma \in \StGraph^{\mathrm{conn}} \\
  g(\Gamma) = g}}
  C[\Gamma]_{\Aut(\Gamma)}.
  \]
  At genus~$0$, only trees survive: $\bar{B}^{(0)}(C) = \Bch(C)$.

\item The genus-$0$ associated graded of the deformation algebra
  is the chiral convolution:
  $\gr_F^0\, L_{\mathrm{mod}}(C) \cong
  \mathfrak{g}^{\mathrm{ch}}_C$.

\item Each $\gr_F^g$ for $g \ge 1$ is a module over
  $\mathfrak{g}^{\mathrm{ch}}_C$ via the adjoint action, and
  the obstruction class
  $\Ob_g(\Theta_{<g}) \in
  H^2(\gr_F^g,\, d_0^{\Theta_0})$
  of Theorem~\textup{\ref{thm:Obg}} lives in the
  cohomology of this module twisted by the genus-$0$ MC element.

\item $[F^p, F^q] \subset F^{p+q}$;
  $D_0(F^p) \subset F^p$;
  $D_1(F^p) \subset F^{p+1}$;
  and $L_{\mathrm{mod}} = \varprojlim_g L_{\mathrm{mod}}/F^g$.

\item The associated spectral sequence has
  $E_1^{p,q} = H^{p+q}(\gr_F^p)$, converges to
  $H^*(L_{\mathrm{mod}})$, and its $d_1$-differential is the
  genus-raising operator~$D_1$.  The zeroth column is
  $H^*(\mathfrak{g}^{\mathrm{ch}}_C)$.
\end{enumerate}
\end{theorem}

\begin{proof}
The differential $D = D_0 + D_1$ respects the genus filtration
because $D_0 = D_{\mathrm{int}} + D_{\mathrm{sep}}$ preserves
genus and $D_1 = D_{\mathrm{nsep}}$ raises it by one
(Proposition~\ref{prop:D0D1}).  The associated graded of the
coalgebra at genus~$g$ picks out stable graphs of total genus
exactly~$g$; at $g = 0$, stability forces $b_1 = 0$ and all
$g(v) = 0$, so only trees remain.

For part~(ii): a coderivation of $\Bmod(C)$ restricts, via
$\Bmod \to \Bmod/F^1 = \Bch$, to a coderivation of the tree-level
coalgebra.  The kernel consists of coderivations supported on
graphs of genus $\ge 1$, which is $F^1\Coder(\Bmod)$.  Thus
$\gr_F^0\,\Coder(\Bmod)[-1] = \Coder(\Bch)[-1] =
\mathfrak{g}^{\mathrm{ch}}_C$.  The bracket on $\gr^0_F$ is the
commutator of tree-level coderivations: at genus~$0$, only
separating one-edge expansions contribute, and their commutator
is again tree-level.

For part~(iii): $[F^0, F^g] \subset F^g$ by the genus-additivity
of separating clutching, so $\gr_F^g$ is a
$\mathfrak{g}^{\mathrm{ch}}$-module.  The twisted differential
$d_0^{\Theta_0}(x) = D_0(x) + [\Theta_0, x]$ is the standard
MC twist.

For parts~(iv)--(v): $[F^p, F^q] \subset F^{p+q}$ because
separating clutching adds genera; completeness is
Theorem~\ref{thm:modular-bar}(iii).  The spectral sequence is
the standard one of a complete filtered dg~Lie algebra.  Since
$D_1 \colon F^p \to F^{p+1}$ and $D_0 \colon F^p \to F^p$, the
induced map $d_1 \colon \gr^p \to \gr^{p+1}$ is $D_1$.
\end{proof}

\begin{remark}[What the genus completion means]
\label{rem:genus-completion-meaning}
\index{genus completion!interpretation}
The theorem says that $L_{\mathrm{mod}}(C)$ is the
pro-nilpotent extension
\[
0 \longrightarrow F^1 L_{\mathrm{mod}}
\longrightarrow L_{\mathrm{mod}}
\longrightarrow \mathfrak{g}^{\mathrm{ch}}_C
\longrightarrow 0
\]
of the chiral convolution by the positive-genus ideal.  The
genus-$0$ chiral convolution $\mathfrak{g}^{\mathrm{ch}}$ controls
tree-level Koszul duality; the modular deformation algebra
$L_{\mathrm{mod}}$ controls the full genus tower.  The passage
from one to the other is governed by the non-separating one-edge
operator $D_1$: it is the algebraic shadow of the first loop
correction, and it is the sole primitive that makes the genus
tower nontrivial.

The corresponding passage on the coalgebra side is equally clean:
\[
\begin{array}{rcl}
\text{trees} & \rightsquigarrow & \text{stable graphs} \\[3pt]
\barB_{\mathrm{op}}(C)
  & \rightsquigarrow
  & \mathrm{FT}_{\cM\!od}(C) \\[3pt]
\mathfrak{g}^{\mathrm{ch}}_C
  & \rightsquigarrow
  & L_{\mathrm{mod}}(C) \\[3pt]
\text{flat twisting } \pi
  & \rightsquigarrow
  & \text{curved twisting } \tau_{\mathrm{mod}} \\[3pt]
d_0^2 = 0
  & \rightsquigarrow
  & d_{\mathrm{fib}}^2 = \kappa \cdot \omega_g
\end{array}
\]
Every entry in the left column is the genus-$0$ associated graded
of the corresponding entry on the right.  Every theorem of
Volume~I at genus~$0$ has a modular extension in the right column.
\end{remark}

\begin{proposition}[Factorisation along boundary strata;
\ClaimStatusProvedHere]
\label{prop:factorisation-boundary-strata}
\index{clutching factorization|textbf}
\index{Maurer--Cartan class!factorisation}
Let $\delta_\cA = D - D^{(0)} = \sum_{g \ge 1} \hbar^g\,
\delta^{(g)}_\cA \in F^1 L_{\mathrm{mod}}$ be the positive-genus
MC element
\textup{(}Theorem~\textup{\ref{thm:coderivation-MC-synthesis})}.
Its genus-$g$ component factorises along the boundary strata
of~$\overline{\cM}_{g,n}$:
\begin{enumerate}[label=\textup{(\roman*)}]
\item \emph{Separating nodes.}
  $\xi_{\mathrm{sep}}^*(\delta^{(g)})
  = \sum_{g_1+g_2=g}
  \delta^{(g_1)} \star \delta^{(g_2)}$,
  where $\star$ is the convolution product through the sewing
  propagator.

\item \emph{Non-separating nodes.}
  $\xi_{\mathrm{ns}}^*(\delta^{(g)})
  = \Delta_{\mathrm{cyc}}(\delta^{(g-1)})
  + \text{lower-genus quadratics}$,
  where $\Delta_{\mathrm{cyc}}$ contracts two inputs with the
  invariant form.

\item \emph{Codimension-two compatibility.}
  On codimension-$2$ strata, the two factorisation maps agree.
  This is the codimension-two cancellation axiom
  (Definition~\ref{def:modular-bar-datum}(3)), transported to
  coderivations.
\end{enumerate}
\end{proposition}

\begin{proof}
Part~(i): $\delta_\cA$ is a coderivation
of~$\Bmod$, so it satisfies the Leibniz rule
$\Delta_{\mathrm{sep}} \circ \delta = (\delta \otimes \id +
\id \otimes \delta) \circ \Delta_{\mathrm{sep}}$
with respect to the separating coproduct.  Projecting to
genus~$g$ and using genus-additivity across the separating edge
gives the factorisation formula.

Part~(ii): the non-separating clutching map pairs two legs on
a genus-$(g{-}1)$ surface; on $\Bmod$, this is
$D_{\mathrm{nsep}}$, whose coderivation incarnation is
$\Delta_{\mathrm{cyc}}$.  The MC equation at genus~$g$
(Theorem~\ref{thm:Obg}) gives
$d_0\delta^{(g)} = -\Delta_{\mathrm{cyc}}(\delta^{(g-1)})
- \tfrac{1}{2}\sum_{h+h'=g} [\delta^{(h)}, \delta^{(h')}]$,
determining $\delta^{(g)}$ from lower genera.

Part~(iii): the codimension-two cancellation axiom says that
for any stable graph with two internal edges, the two orderings
of successive one-edge expansions cancel with signs.  On
coderivations, this is the associativity of iterated
factorisation---equivalently, $\partial^2 = 0$ on the
stratification of~$\overline{\cM}_{g,n}$.
\end{proof}

\begin{remark}[The convolution bracket as correspondence product]
\label{rem:convolution-as-correspondence}
\index{correspondence algebra!modular convolution}
\index{Chriss--Ginzburg!correspondence algebra}
The Lie bracket on $L_{\mathrm{mod}}$ is a correspondence
product.  For $\alpha$ at $(g_1,n_1)$ and $\beta$ at $(g_2,n_2)$,
the bracket $[\alpha,\beta]$ composes through the separating
clutching correspondence
\[
\begin{tikzcd}[column sep=2.5em]
\overline{\cM}_{g_1,n_1+1} \times
\overline{\cM}_{g_2,n_2+1}
& \ar[l, "p"'] Z
  \ar[r, "\xi"] &
\overline{\cM}_{g,n}
\end{tikzcd}
\]
via $[\alpha,\beta] = \xi_* \circ p^*(\alpha \boxtimes \beta)$,
contracted by the invariant pairing on the new edge.  This is
the operadic incarnation of the Chriss--Ginzburg convolution on
Steinberg varieties: as the Steinberg self-intersection
$\widetilde{\cN} \times_{\mathfrak{g}} \widetilde{\cN}$ produces
the Hecke algebra, the modular clutching correspondence produces
the modular convolution.  The Jacobi identity is the associativity
of the clutching correspondence---the commutativity of
codimension-$2$ clutching squares.
\end{remark}

\begin{remark}[Factorisation homology: Ran space to moduli tower]
\label{rem:factoris-ran-to-moduli}
\index{factorisation homology!Ran to moduli tower}
In the language of Beilinson--Drinfeld, the chiral convolution
$\mathfrak{g}^{\mathrm{ch}}_\cA$ is the endomorphism
factorisation homology $\int_{\Ran(X)} \End_\cA$ over the Ran
space of a curve, and the modular convolution $L_{\mathrm{mod}}$
is $\int_{\cM\!od} \End_\cA$---factorisation homology over the
full modular operad $\cM\!od = \{\overline{\cM}_{g,n}\}$.  The
projection $L_{\mathrm{mod}} \to \mathfrak{g}^{\mathrm{ch}}$
is the restriction of factorisation homology along the
inclusion of $\Ran(\P^1)$ (genus-$0$ configurations) into
$\cM\!od$.  The positive-genus MC element $\delta_\cA$ is a
section of the factorisation algebra over the full moduli tower,
and its factorisation along boundary strata
(Proposition~\ref{prop:factorisation-boundary-strata}) is
the factorisation property of this section.

The ribbon-graph variant gives the $E_1$~wing.  Replacing
stable graphs by ribbon graphs---the Feynman transform
$F\!\Ass$ of the associative modular operad---produces the
ordered convolution algebra
$L_{\mathrm{mod}}^{E_1} := \Coder(\Bmod^{\mathrm{rib}})[-1]$.
The coinvariant projection from ribbon to stable graphs is the
passage from colour-ordered to full amplitudes, and the
genus expansion of $L_{\mathrm{mod}}^{E_1}$ is the algebraic
form of the 't~Hooft $1/N$ expansion: planar ribbon graphs
dominate at large~$N$, and the MC element decomposes as
$\sum_g N^{2-2g} \delta_{\mathrm{planar}}^{(g)}$
(Theorem~\ref{thm:formal-genus-expansion}).
\end{remark}

\begin{computation}[The genus completion on the standard families;
\ClaimStatusProvedHere]
\label{comp:genus-completion-standard}
\index{affine!genus completion|textbf}
\begin{enumerate}[label=\textup{(\alph*)}]
\item \emph{Affine} $V^k(\mathfrak{g})$:
  the chiral convolution bracket is
  $[\alpha,\beta] = f^{ab}{}_c\, \alpha \circ_a \beta$.
  The genus-$1$ obstruction is
  $\Ob_1 = (k + h^\vee) \cdot [\omega_1]$.
  The full modular MC class in the scalar sector is the
  $\hat{A}$-genus:
  $\delta_\cA = \sum_{g \ge 1} \hbar^g (k+h^\vee)^g \lambda_g$.

\item \emph{Heisenberg} $\cH_k$:
  $\mathfrak{g}^{\mathrm{ch}} = 0$ (the chiral convolution
  bracket vanishes because the OPE is abelian).  The modular
  deformation is purely genus-raising:
  $D_1$ acts as $\hbar k\, \partial^2/\partial a^i \partial a^j$,
  and the spectral sequence degenerates at~$E_1$.

\item \emph{Virasoro/\/$\cW_N$}:
  all channels---separating bracket, non-separating loop, and
  rigid residues---are active.  The first modular obstruction
  vanishes for $\cW_3$ by the triangular normal-form theorem
  (Theorem~\ref{thm:triangular-vanishing}).
\end{enumerate}
\end{computation}

\subsection{Concrete chain-level presentation from the three foundational papers}
\label{subsec:vol2-three-preprints-chain-level}

\begin{construction}[Homotopy-chiral/log-FM coordinates; \ClaimStatusProvedHere]
\label{const:vol2-three-preprints-chain-level}
Fix an affine cover $\mathfrak U$ of the boundary curve.
The \v{C}ech totalisation
$A^{\mathrm{ch}}_\infty := \check C^\bullet(\mathfrak U;\A)$
is a $\mathrm{Ch}_\infty$-algebra~\cite[Theorem~3.1]{MS24}.
For a logarithmic boundary pair or semistable family,
$\cC^{\log FM}_{X\mathbin{|}D}(n):=C_\bullet(\operatorname{FM}_n(X\mathbin{|}D))$
and $\cC^{\log FM}_{W/B}(n):=C_\bullet(\operatorname{FM}_n(W/B))$.
The modular deformation object in these coordinates is
\[
{\Definfmod}_{\log}(C)
:=
\Convinf\!\Bigl(
\cC^{\log FM}_{\mathrm{mod}},\,
\operatorname{End}_{\mathrm{Ch}_\infty}(A^{\mathrm{ch}}_\infty)
\Bigr),
\]
with strict chart
$L^{\log}_{\mathrm{mod}}(C)
:=
\Convstr(\cC^{\log FM}_{\mathrm{mod}},\,
\operatorname{End}_{\mathrm{Ch}_\infty}(A^{\mathrm{ch}}_\infty))$.
\end{construction}

\begin{remark}[Stable-graph modular bar functor in coordinates]
\label{rem:vol2-log-modular-bar-functor}
The modular bar construction in log-FM coordinates is
\[
\Bmod^{\log}(C)(g,n)
=
\bigoplus_{\Gamma\in \mathsf{Gr}^{\mathrm{st}}_{g,n}}
\left(
\bigotimes_{v\in V(\Gamma)} \bar B^{\mathrm{ch}}(A^{\mathrm{ch}}_\infty;\mathrm{In}(v))
\otimes
C_\bullet(\operatorname{FM}^{\log}_{\Gamma})
\otimes
\orline{E_{\mathrm{int}}(\Gamma)}
\right)_{\Aut(\Gamma)}.
\]
Its differential is the sum of the transferred chiral
differential, separating sewing, planted-forest corrections, and the
non-separating loop operator.
\end{remark}

\begin{computation}[Low-genus control terms; \ClaimStatusProvedHere]
\label{comp:vol2-low-genus-control}
The first coefficients of the modular $L_\infty$-structure in the strict chart:
\[
\ell_3^{(0)}(\alpha,\beta,\gamma)
=
\sum_{T\in\mathsf{Trees}_3}
\frac{1}{|\Aut(T)|}\,
(\pi\circ m_T^{(h)}\circ\iota^{\otimes 3})(\alpha,\beta,\gamma),
\]
\[
\ell_4^{(0)}(\alpha,\beta,\gamma,\delta)
=
\sum_{T\in\mathsf{Trees}_4}
\frac{\pm1}{|\Aut(T)|}\,
(\pi\circ m_T^{(h)}\circ\iota^{\otimes 4})(\alpha,\beta,\gamma,\delta),
\]
and the first genuine modular term is
$\ell_1^{(1)} = \hbar\Delta_{\mathrm{cyc}}$.
For classical $W_3$ the first modular obstruction vanishes.
\end{computation}

\begin{construction}[PVA-coordinate residue formulas; \ClaimStatusProvedHere]
\label{const:vol2-pva-coordinate-graph-coefficients}
For local boundary fields $u_1,\dots,u_n$, the first tree
coefficients are residue integrals over log-FM boundary strata:
\begin{align}
\ell_3^{(0)}(u_1,u_2,u_3)
&=
\sum_{T\in \mathsf{Trees}_3}
\operatorname{Res}_{D_T^{\log}}\,\omega^{T}_{u_1,u_2,u_3},
\label{eq:vol2-ell3-residue}\\
\ell_4^{(0)}(u_1,u_2,u_3,u_4)
&=
\sum_{F\in \mathsf{PF}^{\mathrm{rig}}_4}
\epsilon_F\,\operatorname{Res}_{D_F^{\log}}\,\omega^{F}_{u_1,u_2,u_3,u_4},
\label{eq:vol2-ell4-residue}
\end{align}
with $\ell_1^{(1)}=\hbar\Delta_{\mathrm{cyc}}$.
The three summands of~\eqref{eq:vol2-ell4-residue} are the
contact, additive, and tree rigid types of Mok's degeneration
formula~\cite[Corollary~5.3.4]{Mok25}.
\end{construction}

\begin{computation}[Archetype formulas in boundary coordinates; \ClaimStatusProvedHere]
\label{comp:vol2-archetype-formulas}
Heisenberg ($\{a^i{}_{\lambda}a^j\}=K^{ij}\lambda$):
$\ell_3^{(0)}=0$, $\ell_4^{(0)}=0$,
$\ell_1^{(1)}=\hbar\sum_{i,j}K^{ij}\partial^2/(\partial a^i\partial a^j)$.
Affine currents ($\{J^a{}_{\lambda}J^b\}=f^{ab}{}_cJ^c+k\kappa^{ab}\lambda$):
$[\ell_3^{(0)}]=[f_{abc}\,J^a\otimes J^b\otimes J^c]$,
$\ell_4^{(0)}=0$.
For the $W_3$ PVA, $\Lambda={:}TT{:}-\tfrac{3}{10}\partial^2T$
controls the quartic tree coefficient: $[\ell_4^{(0)}]\propto
\tfrac{16}{22+5c}\,\Lambda$.
\end{computation}

\begin{remark}[Genus-two entry point]
\label{rem:vol2-genus-two-entry}
$\Theta^{(2)}
=
\Theta^{(2)}_{\mathrm{loop}\circ\mathrm{loop}}
+\Theta^{(2)}_{\mathrm{sep}\circ\mathrm{loop}}
+\Theta^{(2)}_{\mathrm{pf}}$.
Heisenberg sees only the first term, affine currents see the first two,
$W_3$/Virasoro exhibit all three.
\end{remark}

\begin{remark}[Rigid-type boundary factorization]
\label{rem:vol2-rigid-type-boundary-factorization}
Mok's degeneration formula writes each rigid boundary component of
$\operatorname{FM}_n(W/B)$ as a push-pull of products of lower log-FM factors:
$\operatorname{FM}_n(W/B)(\rho)\to
\prod_{v\in V(S_\rho)}\operatorname{FM}_{I_v}(Y_v\mathbin{|}D_v)$.
Every clutching identity in the strict modular complex is a
push-pull identity on rigid planted-forest strata.
\end{remark}

%% ----------------------------------------------------------------
%% PRIMITIVE BOUNDARY KERNEL AND BRANCH BV ACTION
%% ----------------------------------------------------------------

\subsection{The primitive boundary kernel, branch BV action,
and the metaplectic square root}
\label{subsec:vol2-primitive-kernel-branch-bv}

The stable-graph expansion of the modular MC field $\Theta_C$ is cofree:
it is uniquely determined by its restriction to primitive corollas and
rigid cutting strata.  That restriction --- the primitive boundary kernel
--- is the native object of the modular PVA theory.  Its finite-rank
branch projection carries a BV master action whose quantum master equation
is the finite-dimensional shadow of the full primitive Maurer--Cartan
equation.

\begin{definition}[Primitive boundary kernel]
\label{def:vol2-primitive-boundary-kernel}
\ClaimStatusProvedHere
\index{primitive boundary kernel|textbf}
In PVA coordinates on the boundary curve $C$, the
\emph{primitive boundary kernel} is
\begin{equation}
\label{eq:vol2-primitive-kernel}
\mathfrak{K}_C
\;:=\;
K_{0,2} + K_{0,3} + K_{0,4}
+ K_{1,1}
+ R_{\mathrm{pf},2} + R_{\mathrm{pf},3} + \cdots,
\end{equation}
where $K_{0,n}$ is the genus-zero $n$-ary corolla coefficient
\textup{(}the $n$-th transferred chiral operation in PVA
coordinates\textup{)}, $K_{1,1}$ is the genus-one/arity-one corolla
\textup{(}the cyclic trace\textup{)}, and $R_{\mathrm{pf},d}$ is the
codimension-$d$ planted-forest residue from Mok's rigid boundary
strata.  This is the PVA-coordinate form of
the primitive logarithmic modular kernel of Volume~I.
\end{definition}

\begin{proposition}[Primitive reconstruction in PVA coordinates]
\label{prop:vol2-primitive-reconstruction}
\ClaimStatusProvedHere
\index{primitive boundary kernel!reconstruction}
The strict modular homological vector field is recovered from
$\mathfrak{K}_C$ by the logarithmic modular Feynman transform:
\begin{equation}
\label{eq:vol2-primitive-reconstruction}
Q_C^{\mathrm{mod}}
\;=\;
(d_{\mathrm{int}} + [\tau,-])
\;+\;
\operatorname{ad}
\operatorname{FT}^{\log}_{\mathrm{mod}}(\mathfrak{K}_C),
\qquad
\Theta_C
\;=\;
\operatorname{FT}^{\log}_{\mathrm{mod}}(\mathfrak{K}_C).
\end{equation}
The fundamental equivalence holds:
$(Q_C^{\mathrm{mod}})^2 = 0
\;\Longleftrightarrow\;
d\,\mathfrak{K}_C + \mathfrak{K}_C \star \mathfrak{K}_C
+ \hbar\,\Delta_{\mathrm{cyc}}(\mathfrak{K}_C) = 0$.
\end{proposition}

\begin{proof}
Coordinate specialization of
the primitive-to-global reconstruction theorem of Volume~I.
\end{proof}

\begin{computation}[Primitive shell equations in boundary coordinates]
\label{comp:vol2-primitive-shell-equations}
\ClaimStatusProvedHere
\index{shell equations!PVA coordinates}
Projecting the primitive master equation to genus two and three gives:
\begin{align}
\label{eq:vol2-primitive-shell-g2}
R_2^{\mathrm{rel}}
&\;=\;
\Delta_{\mathrm{cyc}}(K_{1,1})
\;+\;
\tfrac{1}{2}[K_{1,1},\,K_{1,1}]
\;+\;
\sum_{\rho \in \operatorname{Rig}_2}
\mathfrak{r}_\rho(K_{0,2}),
\\
\label{eq:vol2-primitive-shell-g3}
R_3^{\mathrm{rel}}
&\;=\;
\Delta_{\mathrm{cyc}}(K_{2,\bullet})
\;+\;
[K_{1,1},\,K_{2,\bullet}]
\;+\;
\tfrac{1}{3!}\ell_3(K_{1,1},\,K_{1,1},\,K_{1,1})
\;+\;
\sum_{\rho \in \operatorname{Rig}_2}
\mathfrak{r}_\rho(K_{1,1})
\;+\;
\sum_{\rho \in \operatorname{Rig}_3}
\mathfrak{r}_\rho(K_{0,2}).
\end{align}
The family regimes:
\emph{Heisenberg} sees only $\Delta_{\mathrm{cyc}}$ \textup{(}pure
non-separating loop\textup{)};
\emph{affine currents} add the separating bracket;
\emph{Virasoro}/$\mathcal{W}_N$ activate all channels including rigid
residues.
\end{computation}

\begin{remark}[Notation: $\Delta_{\mathrm{cyc}}$ vs $\Delta_{\mathrm{ns}}$]
\label{rem:vol2-delta-cyc-vs-ns}
The operator $\Delta_{\mathrm{cyc}}$ in \eqref{eq:vol2-primitive-shell-g2}
is the PVA-coordinate form of the non-separating self-gluing operator
$\Delta_{\mathrm{ns}}$ from the primitive shell equations of
Volume~I: both denote the
BV-type contraction that closes one pair of inputs on a single corolla,
producing a genus-raising loop.
\end{remark}

\begin{definition}[Branch BV action in PVA coordinates]
\label{def:vol2-branch-bv-action}
\ClaimStatusProvedHere
\index{branch BV action!PVA coordinates|textbf}
Let $V_C^{\mathrm{br}} = V_C^+ \oplus V_C^-$ be the polarized
branch-active quotient.  The branch BV algebra
$BV_C^{\mathrm{br}} :=
\widehat{\operatorname{Sym}}\bigl((V_C^+)^\vee\bigr)
\widehat\otimes
\Lambda\bigl((V_C^-)^\vee\bigr)$
carries the modular master action
\begin{equation}
\label{eq:vol2-branch-master-action}
S_C^{\mathrm{br}}(x,\xi)
\;:=\;
\tfrac{1}{2}\,\Omega_C(K_C^{\mathrm{br}} x,\, x)
\;+\;
\sum_{m+n \geq 3}
\frac{1}{m!\,n!}\,
\mu^C_{m,n}(x^{\otimes m};\,\xi^{\otimes n}),
\end{equation}
where $\mu^C_{m,n}$ are the branch projections of the primitive
corolla and rigid-cut amplitudes.
\end{definition}

\begin{proposition}[Branch quantum master equation]
\label{prop:vol2-branch-qme}
\ClaimStatusProvedHere
\index{quantum master equation!PVA branch}
The branch differential
$Q_C^{\mathrm{br}} := \{S_C^{\mathrm{br}}, -\}_{\mathrm{BV}}
+ \hbar\,\Delta_C^{\mathrm{br}}$
squares to zero if and only if
\begin{equation}
\label{eq:vol2-branch-qme}
\hbar\,\Delta_C^{\mathrm{br}}\, S_C^{\mathrm{br}}
\;+\;
\tfrac{1}{2}\{S_C^{\mathrm{br}},\, S_C^{\mathrm{br}}\}_{\mathrm{BV}}
\;=\; 0.
\end{equation}
This is the finite-rank BV projection of the primitive
Maurer--Cartan equation for $\mathfrak{K}_C$.
\end{proposition}

\begin{proof}
Coordinate specialization of the branch master equation proved in
Volume~I.
\end{proof}

\begin{computation}[Archetype master actions]
\label{comp:vol2-archetype-master-actions}
\ClaimStatusProvedHere
\index{archetype!master action|textbf}
\index{Heisenberg!primitive kernel}
\index{affine!primitive kernel}
\index{W algebra@$\mathcal{W}_3$!primitive kernel}
The standard families realize distinct primitive kernel regimes:
\begin{enumerate}[label=\textup{(\alph*)}]
\item \emph{Heisenberg} \textup{(}pure quadratic, branch rank $1$\textup{):}
$\mathfrak{K}_{\mathrm{Hei}} = K_{0,2} + K_{1,1}$.
The branch master action is
$S^{\mathrm{br}}_{\mathrm{Hei}}(x) = \tfrac{1}{2}K_{ij}\,x^i x^j$.
\item \emph{Affine $\widehat{\mathfrak{sl}}_2$}
\textup{(}cubic-tree, branch rank $2$\textup{):}
$\mathfrak{K}_{\widehat{\mathfrak{sl}}_2} = K_{0,2} + K_{0,3} + K_{1,1}$.
The branch master action is
$S^{\mathrm{br}}_{\mathrm{aff}}(x) =
\tfrac{1}{2}k\,\kappa_{ab}\,x^a x^b
+ \tfrac{1}{3!}f_{abc}\,x^a x^b x^c$.
\item \emph{$\mathcal{W}_3$}
\textup{(}quartic-rigid, branch rank $3$\textup{):}
$\mathfrak{K}_{\mathcal{W}_3} = K_{0,2} + K_{0,3} + K_{0,4}
+ K_{1,1} + R_{\mathrm{pf},2} + \cdots$.
The quartic quasi-primary
$\Lambda = {:}TT{:} - \tfrac{3}{10}\partial^2 T$ controls the first
genuinely nonlinear term:
$[S^{\mathrm{br}}_{\mathcal{W}_3}]_{\mathrm{quartic}}
\;\propto\;
\dfrac{16}{22 + 5c}\,\Lambda$.
\end{enumerate}
\end{computation}

\begin{corollary}[Metaplectic square root in PVA coordinates]
\label{cor:vol2-metaplectic-square-root}
\ClaimStatusProvedHere
\index{metaplectic half-density!PVA coordinates}
With $T_C^{\mathrm{br,red}}$ the reduced branch transport operator
on $V_C^{\mathrm{br}}$, the formal series
\begin{equation}
\label{eq:vol2-metaplectic}
\delta_C(x)
\;:=\;
\exp\!\Bigl(\tfrac{1}{2}
\operatorname{Tr}\log\bigl(1 - x\,T_C^{\mathrm{br,red}}\bigr)
\Bigr)
\end{equation}
is a canonical square root of the spectral determinant:
$\delta_C(x)^2 = \det(1 - x\,T_C^{\mathrm{br,red}}) = \Delta_C(x)$.
The half-density $\delta_C$ lives on the finite-dimensional polarized
branch sector, not on the full infinite deformation complex.
\end{corollary}

\begin{proof}
Coordinate specialization of
the metaplectic square-root corollary of Volume~I.
\end{proof}

%% ----------------------------------------------------------------
%% END PRIMITIVE BOUNDARY KERNEL
%% ----------------------------------------------------------------

\subsection{Explicit low-order Taylor calculus in PVA coordinates}
\label{subsec:vol2-explicit-low-order-taylor-calculus}

The convolution $L_\infty$ Taylor coefficients~\cite[Theorem~5.1]{RNW19}
are
$\ell_k(f_1,\dots,f_k)(x)
=
\gamma_A(\alpha\circ 1_D)(f_1\otimes\cdots\otimes f_k)\Delta_D^k(x)$.
The secondary Borcherds operations~\cite{MS24} give the mode-level
cubic formula
\begin{equation}
\label{eq:vol2-explicit-l3-secondary-borcherds}
\ell_3^{(0)}(a,b,c)
=
\sum_{p,q,r}
\left(
\int_{\partial^{\log}_{p,q,r}\FM_3} \omega_{p,q,r}
\right)
 j'_{(p,q,r)}(a,b,c),
\end{equation}
where the $j'_{(p,q,r)}$ satisfy the secondary Borcherds identity
\begin{align}
\label{eq:vol2-secondary-borcherds-identity}
&(d\circ j'_{(p,q,r)}+j'_{(p,q,r)}\circ d)(a,b,c)
\notag\\
&\qquad=
\sum_{j\ge 0}\binom{p}{j}
\Bigl(
(-1)^j a_{(p+q-j)}(b_{(r+j)}c)
-(-1)^{j+p+|a||b|}b_{(r+p-j)}(a_{(q+j)}c)
\Bigr)
-\sum_{j\ge 0}\binom{q}{j}(a_{(p+q-j)}b)_{(r+j)}c.
\end{align}

\begin{proposition}[Quartic channels and the first loop term; \ClaimStatusProvedHere]
\label{prop:vol2-quartic-channels-and-loop-term}
The first four Taylor coefficients of the modular deformation field:
\begin{align}
\label{eq:vol2-low-order-atlas}
\Theta_{2}^{(0)} &= H,
\notag\\
\Theta_{3}^{(0)} &= \ell_3^{(0)},
\notag\\
\Theta_{4}^{(0)} &= \Theta_{4}^{\mathrm{contact}} + \Theta_{4}^{12\mid 34} + \Theta_{4}^{13\mid 24} + \Theta_{4}^{14\mid 23},
\notag\\
\Theta_{1}^{(1)} &= \hbar\,\Delta_{\mathrm{cyc}}.
\end{align}
The quartic sector has four primitive graph channels: one zero-edge
contact graph and three one-edge separating graphs.
\end{proposition}

\begin{remark}[Archetype classification]
\label{rem:vol2-family-by-family-low-order-atlas}
\begin{enumerate}[label=\textup{(\roman*)}]
\item Free/abelian: Gaussian ($\ell_3^{(0)}=0$, contact channel vanishes).
\item Affine currents: Lie/tree ($\ell_3^{(0)}\neq 0$, quartic via separating channels only).
\item $\beta\gamma$: contact/quartic ($\ell_3^{(0)}=0$, $\Theta_4^{\mathrm{contact}}\neq 0$).
\item Virasoro/$\mathcal W_N$: mixed (cubic and quartic contact both nonzero).
\end{enumerate}
\end{remark}

\begin{remark}[Rigid-type factorization via Mok's cutting map]
\label{rem:vol2-mok-cutting-map-factorization}
For a rigid combinatorial type $\rho$ of a semistable degeneration, Mok's subdivision
\[
\kappa\colon \Pi_n(\Delta)(\operatorname{cone}(\rho))\to \prod_{v\in V(S_\rho)}\Pi_{I_v}(\Sigma_v)
\]
lifts to the modification
\[
\operatorname{FM}_n(W/B)(\rho)\to \prod_{v\in V(S_\rho)}\operatorname{FM}_{I_v}(Y_v\mathbin{|}D_v).
\]
This is the precise geometric input behind every clutching statement in the PVA quantization chapter:
separating factorization is not an analogy but a push-pull identity on rigid logarithmic boundary pieces.
\end{remark}


\subsection{The genus filtration}

\begin{definition}[Genus-raise filtration]
For $\delta\in L_{\mathrm{mod}}(C)$ define
\[
\delta\in F^pL_{\mathrm{mod}}(C)
\iff
\delta(F^g\Bmod(C))\subseteq F^{g+p}\Bmod(C)
\quad\text{for all }g\ge 0.
\]
\end{definition}

\begin{theorem}[Formal structure theorem; \ClaimStatusProvedHere]\label{thm:formal-structure}
Let $C$ be a modular bar datum. Then:
\begin{enumerate}[leftmargin=2.4em]
\item $L_{\mathrm{mod}}(C)$ is a complete filtered dg Lie algebra with differential
\[
d_L(\delta):=[D,\delta];
\]
\item the filtration is Lie-multiplicative:
\[
[F^pL_{\mathrm{mod}},F^qL_{\mathrm{mod}}]\subseteq F^{p+q}L_{\mathrm{mod}};
\]
\item the positive-genus ideal $F^1L_{\mathrm{mod}}$ is pronilpotent;
\item the genus-zero associated graded piece identifies with the ordinary tree-level deformation complex:
\[
\gr_F^0 L_{\mathrm{mod}}(C)\cong \Coder(\Bch(C))[-1].
\]
\end{enumerate}
\end{theorem}

\begin{proof}
Coderivations of any dg coalgebra form a dg Lie algebra under commutator. Completeness
is inherited from the completeness of $\Bmod(C)$. If $\delta$ raises genus by at least $p$
and $\eta$ by at least $q$, then both compositions $\delta\eta$ and $\eta\delta$ raise genus
by at least $p+q$, proving Lie-multiplicativity. Pronilpotence of $F^1$ follows because on
the finite quotient by $F^N$, sufficiently long brackets vanish.

For the associated graded in genus zero, only total-genus-zero graphs survive, i.e.\ trees
with vertex genera all zero. The nonseparating operator $D_1$ dies in the quotient, so the
differential reduces to the tree-level bar differential.
\end{proof}

\begin{corollary}[No modular theory without loops; \ClaimStatusProvedHere]
If $D_1=0$, then every genus-zero Maurer--Cartan class extends uniquely to a full
Maurer--Cartan class. In particular, the nonseparating one-edge piece is the minimal
algebraic source of genuine modular obstructions.
\end{corollary}

\begin{proof}
Apply \cref{prop:D0D1} and the argument of the proposition from \S3.
\end{proof}

\subsection{Genus-zero comparison with quadratic chiral duality}

\begin{proposition}[Comparison with Gui--Li--Zeng; \ClaimStatusProvedHere]
Suppose the genus-zero truncation of $C(\A)$ is the ordinary quadratic chiral bar datum
of a dualizable quadratic chiral algebra $A_{\cl}(V)$. Then the genus-zero Maurer--Cartan
problem in $\gr_F^0\Defmod(A_{\cl}(V))$ recovers the Gui--Li--Zeng Maurer--Cartan problem:
\[
\Hom(A,B)\hookrightarrow \MC(A^!\otimes B)
\]
in the strictly quadratic case, and the twisted equation
\[
\mu\big((S+\alpha)\boxtimes (S+\alpha)\big)=0
\]
in the inhomogeneous quadratic case.
\end{proposition}

\begin{proof}
On the genus-zero quotient only trees survive, so the modular bar coalgebra reduces to the
ordinary tree-level chiral bar coalgebra. The corresponding coderivation dg Lie algebra is
precisely the disk-level deformation complex. The Gui--Li--Zeng identification is therefore
the genus-zero shadow of the present modular theory.
\end{proof}

\section{Maurer--Cartan obstruction theory}
\label{sec:mc-obstruction-theory}

We now extract the exact obstruction groups.

\subsection{Formal obstruction recursion}

Let $L=L_{\mathrm{mod}}(C)$ and write $d_0=[D_0,-]$, $d_1=[D_1,-]$.

\begin{theorem}[Universal genus-one obstruction; \ClaimStatusProvedHere]\label{thm:Ob1}
Let $\Theta_0\in \MC(\gr_F^0L)$. Then the class
\[
\Ob_1(\Theta_0):=[d_1\Theta_0]
\in
H^2\big((\gr_F^1L)^{\Theta_0},d_0^{\Theta_0}\big),
\qquad
d_0^{\Theta_0}(x):=d_0x+[\Theta_0,x],
\]
is well-defined. Moreover,
\[
\Ob_1(\Theta_0)=0
\iff
\exists\,\Theta_1\in F^1L/F^2L
\text{ such that }
\Theta_0+\Theta_1
\]
solves the Maurer--Cartan equation modulo $F^2L$.
\end{theorem}

\begin{proof}
Expand the Maurer--Cartan equation for $\Theta=\Theta_0+\Theta_1$ modulo $F^2$:
\[
0=d\Theta+\frac12[\Theta,\Theta]
=
\Big(d_0\Theta_0+\frac12[\Theta_0,\Theta_0]\Big)
+
\Big(d_0\Theta_1+d_1\Theta_0+[\Theta_0,\Theta_1]\Big).
\]
The first parenthesis vanishes because $\Theta_0$ is genus-zero Maurer--Cartan, hence the
lifting equation is
\[
d_0^{\Theta_0}\Theta_1=-d_1\Theta_0.
\]
Thus the lift exists iff the cohomology class of $d_1\Theta_0$ vanishes.

To prove that $d_1\Theta_0$ is closed, use \cref{prop:D0D1}. Since $D_0^2=0$,
$D_0D_1+D_1D_0=0$, and $D_1^2=0$, we obtain
\[
d_0^2=0,\qquad d_0d_1+d_1d_0=0.
\]
Applying this to the genus-zero Maurer--Cartan equation
\[
d_0\Theta_0+\frac12[\Theta_0,\Theta_0]=0
\]
and using that $d_1$ is a derivation of the Lie bracket yields
\[
d_0^{\Theta_0}(d_1\Theta_0)=0.
\]
\end{proof}

\begin{theorem}[Higher-genus recursion; \ClaimStatusProvedHere]\label{thm:Obg}
Suppose
\[
\Theta_{<g}:=\Theta_0+\Theta_1+\cdots+\Theta_{g-1}
\]
solves the Maurer--Cartan equation modulo $F^gL$. Then the genus-$g$ forcing term
\[
R_g(\Theta_{<g})
:=
d_1\Theta_{g-1}
+\frac12\sum_{i=1}^{g-1}[\Theta_i,\Theta_{g-i}]
\]
is $d_0^{\Theta_0}$-closed and determines a canonical obstruction class
\[
\Ob_g(\Theta_{<g})
:=
[R_g(\Theta_{<g})]
\in
H^2\big((\gr_F^gL)^{\Theta_0},d_0^{\Theta_0}\big).
\]
It vanishes if and only if there exists $\Theta_g\in F^gL/F^{g+1}L$ such that
\[
\Theta_{\le g}:=\Theta_0+\cdots+\Theta_g
\]
solves the Maurer--Cartan equation modulo $F^{g+1}L$.
\end{theorem}

\begin{proof}
Take the genus-$g$ part of the full Maurer--Cartan equation. Since $D_1$ raises genus by
exactly one, the only differential terms are $d_0\Theta_g$ and $d_1\Theta_{g-1}$, while the
quadratic part is the sum over pairs $(i,g-i)$. Thus
\[
d_0^{\Theta_0}\Theta_g=-R_g(\Theta_{<g}).
\]
The same computation as in \cref{thm:Ob1} shows that $R_g$ is closed. The vanishing
criterion is immediate.
\end{proof}

\begin{corollary}[Full modular lift criterion; \ClaimStatusProvedHere]
A full modular Maurer--Cartan lift
\[
\Theta=\Theta_0+\Theta_1+\Theta_2+\cdots\in \MC(L)
\]
exists if and only if
\[
\Ob_g(\Theta_{<g})=0
\qquad\text{for all }g\ge 1.
\]
\end{corollary}

\subsection{Conditional comparison with the Khan--Zeng half-space theory}

We now formulate the theorem package in the form best suited for later application to the
three-dimensional Poisson sigma model.

\begin{assumption}[Half-space tree-level quantization]
Assume the Khan--Zeng half-space theory on $\mathbb R_+\times X$ with Neumann boundary
condition admits renormalized perturbation theory through tree level.
\end{assumption}

\begin{theorem}[Conditional tree-level boundary quantization; \ClaimStatusConditional]\label{thm:conditional-tree}
Let $V$ be a filtered quadratic PVA and let $\A=\widetilde A_{\cl}(V)$ satisfy the resolved
classical datum hypotheses. Under the half-space tree-level quantization assumption there
exist:
\begin{enumerate}[leftmargin=2.4em]
\item a genus-zero Maurer--Cartan class
\[
\Theta_0\in \MC(\gr_F^0\Defmod(A_{\cl}(V)));
\]
\item a genus-zero boundary chiral/factorization algebra $Q_0(V)$,
\end{enumerate}
such that:
\begin{enumerate}[label=\textnormal{(\alph*)},leftmargin=2.8em]
\item the genus-zero shadow of $\Theta_0$ is the classical $\lambda$-bracket of $V$;
\item the classical limit of $Q_0(V)$ is $A_{\cl}(V)$;
\item under the genus-zero comparison with Gui--Li--Zeng, $\Theta_0$ is the disk-level
Maurer--Cartan class encoding the boundary quantization map
\[
A_{\cl}(V)\longrightarrow Q_0(V).
\]
\end{enumerate}
If $V$ contains a Virasoro element satisfying the Khan--Zeng hypotheses, then $\Theta_0$
lies in the topological subcomplex
\[
\Deftop(A_{\cl}(V))\subseteq \Defmod(A_{\cl}(V)).
\]
\end{theorem}

\begin{remark}
This theorem is conditional because the formal algebraic package constructed here does not
by itself prove the existence of the half-space perturbative quantization. What it does is
identify the \emph{exact deformation-theoretic object} into which that boundary problem
must land.
\end{remark}

\begin{corollary}[Conditional modular quantization criterion; \ClaimStatusConditional]
Under the hypotheses of \cref{thm:conditional-tree}, if all obstruction classes vanish,
then there exists
\[
\Theta\in \MC(\Defmod(A_{\cl}(V)))
\]
whose genus-zero shadow is the classical $\lambda$-bracket and whose boundary value is the
desired quantized chiral algebra. Under the Virasoro hypotheses the lift can be chosen in
$\Deftop(A_{\cl}(V))$.
\end{corollary}

\section{Coordinate realization on local variational polyvectors}

We now construct the concrete coordinate model for the first genus-raising operator.

\subsection{Local polyvectors}

Let
\[
\widehat{\cV}:=
\widehat{\mathbb C}[u^a_{(n)}\mid a=1,\dots,N,\ n\ge 0]
\]
be the completed differential polynomial algebra, with $\partial u^a_{(n)}=u^a_{(n+1)}$.
Introduce odd variables $\theta_a^{(n)}$ satisfying
\[
\partial \theta_a^{(n)}=\theta_a^{(n+1)}.
\]

\begin{definition}[Local polyvectors]
The space of local polyvectors is
\[
\Xloc^\bullet(\widehat{\cV})
:=
\widehat{\cV}[\theta_a^{(n)}]\big/\partial\big(\widehat{\cV}[\theta_a^{(n)}]\big).
\]
We write $\int F$ for the class of $F$ modulo total derivatives. The usual variational
Schouten bracket on $\Xloc^\bullet$ will be denoted by $[\ ,\ ]$.
\end{definition}

\begin{remark}[Shift conventions]
The ordinary Poisson cohomological degree is the polyvector degree. Thus vector fields are
degree $1$ and bivectors are degree $2$. The shifted Maurer--Cartan deformation complex is
obtained from $\Xloc^{\bullet+1}$. In this note we freely pass between the two conventions,
always indicating which degree is meant.
\end{remark}

Every local bivector has a reduced representative of the form
\[
P=\frac12\int \theta_a\,A^{ab}(\partial)\,\theta_b
\]
with $A^{ab}(\partial)$ skew-adjoint and all derivatives moved onto the right-hand factor.

For a filtered quadratic PVA operator
\[
\Pi^{ab}(\partial)=\sum_{m\ge 0}\Pi_m^{ab}\,\partial^m,
\]
write
\[
P_\Pi:=\frac12\int \theta_a\,\Pi^{ab}(\partial)\,\theta_b.
\]
The PVA Jacobi identity is equivalent to the variational master equation
\[
[P_\Pi,P_\Pi]=0.
\]

\subsection{The reduced odd Laplacian}

\begin{definition}[Reduced cyclic odd Laplacian]
Define
\[
\Delta_{\mathrm{cyc}}\!\left(\int F\right)
:=
\sum_{a=1}^N\sum_{n\ge 0}
\int
(-\partial)^n
\frac{\partial^2 F}{\partial u^a_{(n)}\,\partial \theta_a}.
\]
\end{definition}

\begin{remark}[Interpretation]
The operator $\Delta_{\mathrm{cyc}}$ is the coordinate model of the nonseparating one-edge
expansion. Geometrically it contracts the two half-edges of a one-loop self-gluing; in BV
language it is the reduced divergence of a local bivector.
\end{remark}

\begin{theorem}[Coordinate formula for the genus-raising operator; \ClaimStatusProvedHere]\label{thm:D1formula}
Let $\Pi=(\Pi^{ab}(\partial))$ be a filtered quadratic PVA operator. Then
\[
D_1(P_\Pi)=\Delta_{\mathrm{cyc}}P_\Pi
=
\int \sum_{b=1}^N \mathfrak M_\Pi^b(\partial)(\theta_b),
\]
where
\[
\mathfrak M_\Pi^b(\partial)
:=
\sum_{a=1}^N \frac{\delta \Pi^{ab}(\partial)}{\delta u^a},
\qquad
\frac{\delta \Pi^{ab}(\partial)}{\delta u^a}
:=
\sum_{n\ge 0}(-\partial)^n\circ
\frac{\partial \Pi^{ab}(\partial)}{\partial u^a_{(n)}}.
\]
We call $\mathfrak M_\Pi=(\mathfrak M_\Pi^b)_b$ the \emph{variational modular operator}
of $\Pi$.
\end{theorem}

\begin{proof}
Differentiate
\[
P_\Pi=\frac12\int \theta_a\Pi^{ab}(\partial)\theta_b
\]
with respect to $\theta_a$ to contract one odd half-edge, then take the variational
derivative with respect to the even jet variables $u^a_{(n)}$. Each integration by parts
contributes a factor $(-\partial)^n$. Summing over $a$ and $n$ produces exactly the
displayed formula.
\end{proof}

\begin{lemma}[BV properties of $\Delta_{\mathrm{cyc}}$; \ClaimStatusProvedHere]\label{lem:BV}
The operator $\Delta_{\mathrm{cyc}}$ is odd, second-order, and satisfies:
\begin{align*}
\Delta_{\mathrm{cyc}}^2&=0,\\
\Delta_{\mathrm{cyc}}[X,Y]
&=
[\Delta_{\mathrm{cyc}}X,Y]
+(-1)^{|X|+1}[X,\Delta_{\mathrm{cyc}}Y].
\end{align*}
\end{lemma}

\begin{proof}
This is the standard BV-coordinate computation in jet coordinates. The Euler-operator
factors $(-\partial)^n$ already incorporate all integration-by-parts signs. Super-commutativity
forces the square to vanish, and the second-order nature yields the derivation rule for the
Schouten bracket.
\end{proof}

\subsection{The variational modular class}

\begin{theorem}[Variational modular class; \ClaimStatusProvedHere]\label{thm:modclass}
Let $\Pi$ be a filtered quadratic PVA operator. Then
\[
[P_\Pi,\Delta_{\mathrm{cyc}}P_\Pi]=0.
\]
Hence $\Delta_{\mathrm{cyc}}P_\Pi$ determines a canonical local Poisson cohomology class
\[
\Mod(V):=[\Delta_{\mathrm{cyc}}P_\Pi]\in
H^1_{\delta_\Pi}\!\big(\Xloc^\bullet(\widehat{\cV})\big),
\qquad
\delta_\Pi:=[P_\Pi,-].
\]
We call $\Mod(V)$ the \emph{variational modular class} of the filtered quadratic PVA.
\end{theorem}

\begin{proof}
Apply \cref{lem:BV} to the identity $[P_\Pi,P_\Pi]=0$:
\[
0=\Delta_{\mathrm{cyc}}[P_\Pi,P_\Pi]=2[P_\Pi,\Delta_{\mathrm{cyc}}P_\Pi].
\]
Thus $\Delta_{\mathrm{cyc}}P_\Pi$ is $\delta_\Pi$-closed.
\end{proof}

\begin{remark}[Obstruction shadow]
The class $\Mod(V)$ lives in ordinary Poisson cohomology. Under the one-edge comparison
between the nonseparating bar operator and the local variational model, it is the coordinate
shadow of the genus-one forcing term $d_1\Theta_0$ in \cref{thm:Ob1}. In particular, if
$\Delta_{\mathrm{cyc}}P_\Pi=0$ then the coordinate shadow of the first modular obstruction
vanishes.
\end{remark}

\begin{lemma}[Gauge-independence up to exact terms; \ClaimStatusProvedHere]\label{lem:gaugeind}
Let $f$ be an even local functional and define
\[
\Delta_f:=e^{-f}\Delta_{\mathrm{cyc}}e^f.
\]
Then
\[
\Delta_fP_\Pi-\Delta_{\mathrm{cyc}}P_\Pi=[P_\Pi,f].
\]
Hence the cohomology class $\Mod(V)$ is independent of the choice of local density
or renormalization scheme.
\end{lemma}

\begin{proof}
Expand $e^{-f}\Delta_{\mathrm{cyc}}e^f$ on the bivector $P_\Pi$. Since $f$ has polyvector
degree zero, the new term is precisely the Hamiltonian bracket with $f$.
\end{proof}

\section{A vanishing theorem for triangular $W$-normal form}

We now identify a broad class of $W$-type PVAs for which the first variational modular
class vanishes identically.

\begin{definition}[Triangular $W$-normal form]
A filtered quadratic PVA is in \emph{triangular $W$-normal form} if it is strongly
generated by
\[
T,\ W^1,\dots,W^r
\]
such that:
\begin{enumerate}[leftmargin=2.4em]
\item $T$ is Virasoro:
\[
\{T_\lambda T\}=\partial T+2T\lambda+K(\lambda),
\]
with $K(\lambda)$ independent of the fields;
\item each $W^\alpha$ is primary of weight $h_\alpha$:
\[
\{T_\lambda W^\alpha\}=\partial W^\alpha+h_\alpha W^\alpha\lambda;
\]
\item in the row indexed by $W^\alpha$, the only dependence on jets of $W^\alpha$
appears in the entry $\Pi^{\alpha T}$.
\end{enumerate}
\end{definition}

\begin{theorem}[Triangular vanishing; \ClaimStatusProvedHere]\label{thm:triangular-vanishing}
If a filtered quadratic PVA is in triangular $W$-normal form, then
\[
\Delta_{\mathrm{cyc}}P_\Pi=(r+1)\int \partial\theta_T=0
\qquad\text{in }\Xloc^1/\partial.
\]
Consequently,
\[
\Mod(V)=0.
\]
\end{theorem}

\begin{proof}
The Virasoro row has operator
\[
\Pi^{TT}=T'+2T\partial+K(\partial),
\]
so
\[
\frac{\delta \Pi^{TT}}{\delta T}=(-\partial)\circ 1 + 2\partial=\partial.
\]

For a primary field $W^\alpha$ of weight $h_\alpha$,
\[
\Pi^{T\alpha}=(W^\alpha)'+h_\alpha W^\alpha\partial.
\]
Skew-adjointness gives
\[
\Pi^{\alpha T}=(h_\alpha-1)(W^\alpha)'+h_\alpha W^\alpha\partial,
\]
hence
\[
\frac{\delta \Pi^{\alpha T}}{\delta W^\alpha}
=-(h_\alpha-1)\partial+h_\alpha\partial=\partial.
\]
By assumption there are no further appearances of $W^\alpha$ in the $\alpha$-row. Therefore
the total modular operator contributes $\partial$ once from the Virasoro row and once from
each primary row:
\[
\Delta_{\mathrm{cyc}}P_\Pi=(r+1)\int \partial \theta_T.
\]
This vanishes modulo total derivatives.
\end{proof}

\begin{remark}
This theorem gives a large natural vanishing class. It says that in a broad swath of
classical $W$-type examples the first modular anomaly does \emph{not} appear as an
obstruction class. The quantum problem starts instead with the torsor of genus-one lifts.
\end{remark}

\section{The holographic genus expansion}
\label{sec:holographic-genus-expansion}

The modular bar complex carries a genus expansion that formally
resembles the 't~Hooft $1/N$ expansion of gauge theories.  We
make the formal structure precise and separate what is proved from
what requires new input.

\begin{theorem}[Formal genus expansion of the modular bar complex]
\label{thm:formal-genus-expansion}
Let $C$ be a modular bar datum.  The total differential of the
modular bar coalgebra $\Bmod(C)$ admits the genus expansion
\begin{equation}\label{eq:genus-expansion}
D \;=\; \sum_{g \ge 0} \hbar^g\, D^{(g)},
\end{equation}
where $D^{(g)}$ is the contribution from stable graphs of total
genus~$g$:
\begin{equation}\label{eq:Dg-graph-sum}
D^{(g)} \;=\; \sum_{\substack{\Gamma \in
\StGraph^{\mathrm{conn}} \\ g(\Gamma) = g}}
\frac{1}{|\Aut(\Gamma)|}\, D_\Gamma.
\end{equation}
The identity $D^2 = 0$ decomposes by genus into
\begin{equation}\label{eq:D-squared-genus}
\sum_{p+q=g} D^{(p)} D^{(q)} = 0
\qquad (g \ge 0).
\end{equation}
In particular:
\begin{enumerate}[label=\textup{(\roman*)}]
\item \textup{[\ClaimStatusProvedHere]} At genus $0$: $(D^{(0)})^2 = 0$
  \textup{(}the tree-level bar differential squares to zero\textup{)}.
\item \textup{[\ClaimStatusProvedHere]} At genus $1$:
  $D^{(0)} D^{(1)} + D^{(1)} D^{(0)} = 0$
  \textup{(}the one-loop correction is a $D^{(0)}$-derivation\textup{)}.
\item \textup{[\ClaimStatusProvedElsewhere]} The full expansion is the Feynman transform of the modular operad
  $\cM\!od = \{\overline{\cM}_{g,n}\}$\textup{:}
  $\bar B^{\mathrm{full}}(\cA) \simeq \mathrm{FT}_{\cM\!od}(\cA)$
  by the higher-genus prism principle of Volume~I.
\end{enumerate}
\end{theorem}

\begin{proof}
Equation~\eqref{eq:genus-expansion} is the definition: the genus
filtration on $\Bmod(C)$ induces a grading on the differential.
The graph sum~\eqref{eq:Dg-graph-sum} counts each isomorphism class
of stable graph once, with its automorphism weight.
The genus decomposition of $D^2 = 0$
(equation~\eqref{eq:D-squared-genus}) follows from
Theorem~\ref{thm:modular-bar} and the genus-additivity of
composition: the composition of a genus-$p$ graph with a genus-$q$
graph produces a genus-$(p+q)$ graph.  Parts~(i) and~(ii)
specialise the genus decomposition.  Part~(iii) is the
higher-genus prism principle of Volume~I.
\end{proof}

\begin{remark}[The formal 't~Hooft structure]
\label{rem:thooft-formal}
\index{t Hooft expansion@'t~Hooft expansion}
The genus expansion~\eqref{eq:genus-expansion} has the same
combinatorial structure as the 't~Hooft $1/N$ expansion of a
gauge theory with gauge group $\mathrm{GL}_N$:
\begin{itemize}
\item Contributions are indexed by stable graphs (= ribbon graphs
  in the 't~Hooft language).
\item Each graph $\Gamma$ carries the symmetry factor
  $1/|\Aut(\Gamma)|$.
\item The genus parameter $\hbar$ organises the expansion.
\item The string-coupling weight $\hbar^{2g(\Gamma)-2}$ is the
  Euler characteristic weight.
\end{itemize}
The identification $\hbar \leftrightarrow g_s^2 \leftrightarrow 1/N$
would be the holographic identification.
\end{remark}

\begin{remark}[What is proved and what is not]
\label{rem:thooft-not-proved}
The following are \textbf{proved}:
\begin{enumerate}[leftmargin=2.4em]
\item The genus expansion exists and $D^2 = 0$ is equivalent to
  the genus-by-genus identities~\eqref{eq:D-squared-genus}.
\item The Maurer--Cartan element $\delta_\cA = D - D^{(0)}$
  lives in the positive-genus filtration and its MC equation is the
  obstruction tower for the modular lift
  (Theorem~\ref{thm:coderivation-MC-synthesis}).
\item The expansion is the Feynman transform of the modular operad.
\end{enumerate}
The following are \textbf{not proved in either volume}:
\begin{enumerate}[leftmargin=2.4em]
\item The identification of $\hbar$ with a large-$N$ parameter for
  any specific family of theories.
\item The existence of a family $\cA_N$ of logarithmic
  $\SCchtop$-algebras parametrised by $N$ whose large-$N$ limit
  recovers the genus expansion.
\item Bulk reconstruction from the boundary datum.
\end{enumerate}
\end{remark}

\begin{conjecture}[Holographic genus identification for $\mathrm{GL}_N$ Chern--Simons]%
\label{conj:thooft-gl-n}
\ClaimStatusConjectured.
\index{t Hooft expansion@'t~Hooft expansion!holographic identification}
Let $\cA_N$ denote the boundary vertex algebra of the half-space BV
quantization of the $\mathrm{GL}_N$ Chern--Simons theory on
$\C \times \R_{\ge 0}$, i.e.\ the level-shifted affine Kac--Moody
algebra $V^{K_{\mathrm{eff}}}(\mathfrak{gl}_N)$.  Then:
\begin{enumerate}[label=\textup{(\alph*)}]
\item The genus-$g$ Maurer--Cartan coefficient
  $\delta_{\cA_N}^{(g)}$ admits a $1/N$ expansion:
  \[
  \delta_{\cA_N}^{(g)} = N^{2-2g}\, \delta_{\mathrm{planar}}^{(g)}
  + O(N^{1-2g}).
  \]
\item In the large-$N$ limit, the planar part
  $\delta_{\mathrm{planar}}^{(0)}$ encodes the classical
  chiral algebra, and the genus-$g$ corrections
  $\delta_{\mathrm{planar}}^{(g)}$ encode the $g$-loop
  corrections to the boundary OPE.
\item The genus expansion of the modular bar complex, with
  $\hbar = 1/N$, recovers the $1/N$ expansion of the boundary
  correlators.
\end{enumerate}
The holographic identification is then: the modular bar complex at
genus $g$ computes the $1/N^{2g}$ correction to the boundary
theory of $N$ branes.
\end{conjecture}

\begin{remark}[Why $\mathrm{GL}_N$ and not $\mathrm{SL}_N$]
The $\mathrm{GL}_N$ theory has $U(1)$ center that decouples in the
large-$N$ limit, making the 't~Hooft counting clean: each
ribbon graph contributes $N^{\chi(\Gamma)}$, where
$\chi(\Gamma) = 2 - 2g(\Gamma)$ is the Euler characteristic.
For $\mathrm{SL}_N$, the center is $\mathbb{Z}/N$ and the counting is more
subtle.
\end{remark}
